\subsection{Detector Comparison and Reporting Scope}
\label{sec:detector_results_scope}

We propose \textbf{SAFR-YOLO} as the 2D evidence generator for the full
CoM3D-ACE system and report it as \textbf{Ours} in all comparison tables.
P2P4-SelfAttnFR denotes the selected implementation and ablation component
combination. The detector contribution is evaluated with two
paper-facing views: a normalized three-seed comparison against strong
YOLO-family, transformer, and cited related-work baselines, plus a compact
ablation table for the selected modules. This keeps the main paper
self-contained while preventing exploratory single-seed architecture-search
rows from being mixed with the final detector claim.

\paragraph{Implementation and evaluation setup.}
Unless otherwise specified, detector experiments use VisDrone2019-DET with
1280-pixel input resolution and the official validation split. The converted
YOLO-format benchmark keeps the 10 VisDrone categories: pedestrian, people,
bicycle, car, van, truck, tricycle, awning-tricycle, bus, and motor. The
paper-facing three-seed protocol uses seeds 42, 123, and 2026; each row reports
the mean and standard deviation over the completed seeds. Training uses batch
size 4 for large/proposed 1280-pixel runs, four dataloader workers, automatic
optimizer selection, AMP, deterministic mode, and the best validation
checkpoint. We do not add custom paper-specific augmentation; all promoted rows
share the exported Ultralytics train-time recipe in their \texttt{args.yaml}
files, including standard HSV, translation/scale, horizontal flip, mosaic, and
close-mosaic settings. The training server contains an Intel Xeon Silver 4510
CPU with 12 physical cores, 24 threads, 125 GiB system memory, and two RTX 5090
GPUs. Parameter counts and GFLOPs are taken from the corresponding model
summaries under the same input protocol. Remaining queue runs use shorter
early-stopping patience only for screening; final paper rows are promoted after
the normalized three-seed record is complete.

\paragraph{Main comparison rule.}
The main normalized detector table contains the selected detector, the strongest
completed YOLO-family and transformer baselines, and the completed cited
related-work reproductions under the same VisDrone2019-DET 1280-pixel,
three-seed protocol. Detailed related-work reproduction status and protocol
caveats are reported in the supplementary material. The full YOLO
nano/small/medium/large scale sweep, protocol details, heat maps, NMS
sensitivity, and module-search trails are also reported there. The main table
therefore focuses on the rows that directly support the detector claim under a
single controlled protocol.
At the current checkpoint, six cited related-work detector rows have normalized
1280-pixel three-seed records: CSFPR-RTDETR, MFFSODNet, SFFEF-YOLO, BPD-YOLO,
HF-D-FINE, and UAVDet.

\begin{table*}[t]
\centering
\scriptsize
\setlength{\tabcolsep}{3pt}
\caption{Main detector comparison on VisDrone validation at 1280 input resolution. All rows are completed three-seed runs under the same protocol; our proposed detector is placed at the bottom following common comparison-table style.}
\label{tab:detector_preview_overall}
\resizebox{\textwidth}{!}{%
\begin{tabular}{@{}lrrrrrrrr@{}}
\toprule
Method & AP & AP50 & P & R & F1 & Params(M) & GFLOPs & $\Delta$AP \\
\midrule
\multicolumn{9}{@{}l}{\textit{Cited UAV / related-work detector reproductions}} \\
UAVDet [16] & 0.3812 $\pm$ 0.0013 & 0.6005 $\pm$ 0.0036 & 0.6542 & 0.5899 & 0.6203 & 33.55 & 201.20 & +0.0035 \\
BPD-YOLO [7] & 0.3804 $\pm$ 0.0035 & 0.5986 $\pm$ 0.0027 & 0.6626 & 0.5802 & 0.6187 & 24.53 & 109.50 & +0.0028 \\
SFFEF-YOLO [4] & 0.3749 $\pm$ 0.0040 & 0.5914 $\pm$ 0.0039 & 0.6540 & 0.5781 & 0.6137 & 20.87 & 108.70 & -0.0027 \\
\midrule
\multicolumn{9}{@{}l}{\textit{YOLO-family baselines}} \\
YOLOv11l & 0.3777 $\pm$ 0.0004 & 0.5981 $\pm$ 0.0011 & 0.6666 & 0.5881 & 0.6248 & 25.32 & 87.30 & +0.0000 \\
YOLOv12l & 0.3771 $\pm$ 0.0020 & 0.5958 $\pm$ 0.0025 & 0.6698 & 0.5805 & 0.6219 & 26.40 & 89.40 & -0.0006 \\
YOLOv8l & 0.3765 $\pm$ 0.0003 & 0.5963 $\pm$ 0.0008 & 0.6695 & 0.5770 & 0.6198 & 43.64 & 165.40 & -0.0011 \\
YOLOv9c & 0.3738 $\pm$ 0.0009 & 0.5942 $\pm$ 0.0020 & 0.6641 & 0.5831 & 0.6210 & 25.54 & 103.70 & -0.0039 \\
YOLOv9m & 0.3700 $\pm$ 0.0003 & 0.5910 $\pm$ 0.0012 & 0.6708 & 0.5749 & 0.6192 & 20.17 & 77.60 & -0.0077 \\
\midrule
\textbf{Ours} & \textbf{0.3822 $\pm$ 0.0007} & \textbf{0.6052 $\pm$ 0.0012} & \textbf{0.6732} & \textbf{0.5872} & \textbf{0.6273} & \textbf{20.82} & \textbf{109.30} & \textbf{+0.0045} \\
\bottomrule
\end{tabular}
}
\end{table*}


\paragraph{Main detector result analysis.}
Table~\ref{tab:detector_preview_overall} shows that the final detector gives
the highest AP and AP50 among the completed 1280-pixel three-seed rows. The
margin over YOLOv11l is modest but consistent across seeds, and the gain is
obtained with 20.82M parameters rather than the 25.32M parameters of YOLOv11l
or the 33.55M parameters of the strongest reproduced related-work row,
UAVDet~\cite{yang2026uavdet}. The comparison also clarifies the trade-off:
SAFR-YOLO is not the lowest-FLOP model, but it occupies the best
accuracy--parameter point among the paper-facing rows, which is the relevant
constraint for our cooperative UAV evidence generator.

\paragraph{Related-work comparison and coverage.}
Among cited UAV small-object detector reproductions, UAVDet~\cite{yang2026uavdet} is the
strongest completed row under the normalized 1280-pixel, three-seed protocol.
Its reproduced mean AP/AP50 is 0.3812/0.6005 with 33.55M parameters. Ours
remains slightly higher at 0.3822/0.6052 while using 20.82M parameters. We
therefore report UAVDet as an inspired reproduction in the supplementary
protocol table rather than treating it as an official checkpoint result.

\paragraph{Selected detector.}
The current selected detector, SAFR-YOLO, is reported as \textbf{Ours} in
result tables. Its implementation keeps a high-resolution P2/P3/P4 detection
path and adds SelfAttnFR refinement with TinyFReLU; this implementation is
abbreviated as P2P4-SelfAttnFR in ablation and engineering artifacts. Compared
with YOLOv11l under the same three-seed protocol, it
improves AP from 0.3777 to 0.3822 and AP50 from 0.5981 to 0.6052 while reducing
the parameter count from 25.32M to 20.82M. This is the clean main-paper claim:
better AP/AP50 and fewer parameters than the strongest normalized stock YOLO
baseline. The trade-off cost is higher GFLOPs, so we do not claim a
speed or compute reduction unless later profiling supports it.

The paired-seed test against YOLOv11l gives AP $p=0.0126$ and AP50 $p=0.0223$
with seeds 42, 123, and 2026. Precision, recall, and F1 remain secondary
supporting metrics because their paired significance is weaker at the current
sample size.

We show the detector trade-off through AP/AP50, parameter count, GFLOPs, and
the AP--parameter figure instead of adding a single hand-designed trade-off
score to the main table. This keeps the table interpretable and avoids making
the conclusion depend on an arbitrary weighting between accuracy and model
cost.

\paragraph{Paper-ready detector figures.}
The main paper reserves space for one compact quantitative detector figure and
one qualitative detector figure after final rendering. The compact quantitative
figure visualizes AP/AP50 versus parameter count for
Table~\ref{tab:detector_preview_overall}, while the qualitative figure compares
YOLOv11l and Ours on crowded UAV crops with adjacent objects. The
Grad-CAM-style feature activation heat maps for YOLOv11l, YOLOv9c, and
Ours are provided in the supplementary material as an additional qualitative
explanation of where the detectors respond on crowded small-object crops.

\paragraph{Ablation reporting.}
The main paper reports the completed compact final ablation in
Table~\ref{tab:main_detector_ablation}. The table isolates the contribution of
the high-resolution P2/P3/P4 heads, TinyFReLU, and SelfAttn refinement before
the full SelfAttnFR detector. The final detector gives the strongest AP, AP50,
and F1 among these interpretable variants. Single-seed search candidates such as
DCT, wavelet, SE, DynFreq, CBAM, and other module probes are kept in the
supplementary material as design evidence unless they become part of the
selected detector.

\begin{table*}[t]
\centering
\scriptsize
\setlength{\tabcolsep}{3pt}
\caption{Final detector ablation under the normalized VisDrone2019-DET
1280-pixel protocol. Deltas are relative to YOLOv11l. AP and AP50 are reported
as mean $\pm$ standard deviation over seeds 42, 123, and 2026 when available.}
\label{tab:main_detector_ablation}
\begin{tabular}{lcccccc}
\toprule
Component & AP & AP50 & F1 & $\Delta$AP & Params & Seeds \\
\midrule
YOLOv11l baseline & 0.3777 $\pm$ 0.0004 & 0.5981 $\pm$ 0.0011 & 0.6248 & +0.0000 & 25.32M & 3 \\
P2/P3/P4 heads & 0.3807 $\pm$ 0.0020 & 0.6028 $\pm$ 0.0043 & 0.6248 $\pm$ 0.0042 & +0.0030 & 20.82M & 3 \\
P2/P3/P4 + TinyFReLU & 0.3784 $\pm$ 0.0009 & 0.5987 $\pm$ 0.0018 & 0.6199 $\pm$ 0.0028 & +0.0008 & 20.82M & 3 \\
P2/P3/P4 + SelfAttn & 0.3815 $\pm$ 0.0004 & 0.6038 $\pm$ 0.0011 & 0.6256 $\pm$ 0.0025 & +0.0038 & 20.82M & 3 \\
P2/P3/P4 + SelfAttnFR & 0.3822 $\pm$ 0.0007 & 0.6052 $\pm$ 0.0012 & 0.6273 & +0.0045 & 20.82M & 3 \\
\bottomrule
\end{tabular}

\end{table*}

\paragraph{Ablation analysis.}
Table~\ref{tab:main_detector_ablation} indicates that most of the detector gain
comes from preserving high-resolution P2/P3/P4 prediction heads, which improves
AP by 0.0030 over YOLOv11l while reducing parameters. Self-attention refinement
adds a further AP/AP50 gain, and the full SelfAttnFR variant gives the best
combined AP, AP50, and F1. TinyFReLU alone is not sufficient to improve the
full detector, but it remains part of the final block because it supports the
local frequency-refinement path when coupled with SelfAttn. This pattern
supports the proposed design as a compact combination of high-resolution heads,
context refinement, and spatial-frequency gating rather than a single isolated
module.
